\section{Canonical Quantization}

\subsection{Klein--Gordon Equation (Classical)}
For a real scalar field:
\begin{equation}
\mathcal{L}=\frac12 \partial_\mu\phi\,\partial^\mu\phi-\frac12 m^2\phi^2
\qquad\Rightarrow\qquad
(\square+m^2)\phi=0.
\end{equation}
Plane-wave solutions:
\begin{equation}
\phi(x)\sim e^{-ip\cdot x},\qquad p^2=m^2,\qquad p^0=\pm E_{\mathbf{p}},\quad
E_{\mathbf{p}}=\sqrt{\mathbf{p}^2+m^2}.
\end{equation}

\subsection{Warm-up: Simple Harmonic Oscillator (SHO)}
Hamiltonian:
\begin{equation}
H=\frac{p^2}{2m}+\frac12 m\omega^2 x^2.
\end{equation}
Define ladder operators:
\begin{equation}
a=\sqrt{\frac{m\omega}{2}}\left(x+\frac{i}{m\omega}p\right),\qquad
a^\dagger=\sqrt{\frac{m\omega}{2}}\left(x-\frac{i}{m\omega}p\right),
\end{equation}
with
\begin{equation}
[a,a^\dagger]=1.
\end{equation}
Then
\begin{equation}
H=\omega\left(a^\dagger a+\frac12\right),
\qquad
E_n=\omega\left(n+\frac12\right),\quad n=0,1,2,\dots
\end{equation}

\subsection{Canonical Quantization of a Real Scalar Field}
Promote $\phi,\pi$ to operators and impose equal-time commutators:
\begin{equation}
\pi(x)=\frac{\partial\mathcal{L}}{\partial\dot{\phi}}=\dot{\phi},
\end{equation}
\begin{equation}
[\phi(\mathbf{x},t),\pi(\mathbf{y},t)]=i\,\delta^{(3)}(\mathbf{x}-\mathbf{y}),
\qquad
[\phi(\mathbf{x},t),\phi(\mathbf{y},t)]=0,
\qquad
[\pi(\mathbf{x},t),\pi(\mathbf{y},t)]=0.
\end{equation}

\subsection{Mode Expansion and Free Quantum Field}
The general operator solution of the KG equation is expanded as:
\begin{equation}
\phi(x)=\int \frac{d^3p}{(2\pi)^3}\,\frac{1}{\sqrt{2E_{\mathbf{p}}}}
\left(
a_{\mathbf{p}}e^{-ip\cdot x}+a^\dagger_{\mathbf{p}}e^{ip\cdot x}
\right),
\qquad p^0=E_{\mathbf{p}}.
\end{equation}
Imposing canonical commutators implies:
\begin{equation}
[a_{\mathbf{p}},a^\dagger_{\mathbf{q}}]=(2\pi)^3\delta^{(3)}(\mathbf{p}-\mathbf{q}),
\qquad
[a_{\mathbf{p}},a_{\mathbf{q}}]=0,
\qquad
[a^\dagger_{\mathbf{p}},a^\dagger_{\mathbf{q}}]=0.
\end{equation}

\subsection{Hamiltonian and Vacuum Energy}
For the free KG field,
\begin{equation}
\mathcal{H}=\frac12\pi^2+\frac12(\nabla\phi)^2+\frac12 m^2\phi^2,
\qquad
H=\int d^3x\,\mathcal{H}.
\end{equation}
In terms of creation/annihilation operators:
\begin{equation}
H=\int \frac{d^3p}{(2\pi)^3}\,E_{\mathbf{p}}
\left(a^\dagger_{\mathbf{p}}a_{\mathbf{p}}+\frac12(2\pi)^3\delta^{(3)}(0)\right).
\end{equation}
The $\delta^{(3)}(0)$ term corresponds to the (divergent) \textbf{vacuum energy}:
\begin{equation}
E_0=\frac12\sum_{\mathbf{p}}E_{\mathbf{p}} \;\;\to\;\; \infty.
\end{equation}

\paragraph{Normal ordering.}
Define normal ordering by moving all creation operators to the left:
\begin{equation}
:H: \;=\int \frac{d^3p}{(2\pi)^3}\,E_{\mathbf{p}}\,a^\dagger_{\mathbf{p}}a_{\mathbf{p}},
\end{equation}
so that
\begin{equation}
\langle 0|:H:|0\rangle=0.
\end{equation}

\subsection{Vacuum and Particles}
Define the vacuum by
\begin{equation}
a_{\mathbf{p}}|0\rangle=0 \quad \forall \mathbf{p}.
\end{equation}
One-particle state:
\begin{equation}
|\mathbf{p}\rangle = a^\dagger_{\mathbf{p}}|0\rangle.
\end{equation}
Multi-particle states:
\begin{equation}
|\mathbf{p}_1,\dots,\mathbf{p}_n\rangle
=a^\dagger_{\mathbf{p}_1}\cdots a^\dagger_{\mathbf{p}_n}|0\rangle.
\end{equation}
Number operator:
\begin{equation}
N=\int \frac{d^3p}{(2\pi)^3}\,a^\dagger_{\mathbf{p}}a_{\mathbf{p}}.
\end{equation}

\subsection{Relativistic Normalization}
A common convention is
\begin{equation}
\langle \mathbf{p}|\mathbf{q}\rangle=(2\pi)^3\delta^{(3)}(\mathbf{p}-\mathbf{q}).
\end{equation}
A Lorentz-invariant measure is
\begin{equation}
\int \frac{d^3p}{(2\pi)^3}\,\frac{1}{2E_{\mathbf{p}}}.
\end{equation}
Covariant normalization often used in scattering:
\begin{equation}
|p\rangle_{\mathrm{cov}}=\sqrt{2E_{\mathbf{p}}}\,|\mathbf{p}\rangle,
\qquad
\langle p|q\rangle_{\mathrm{cov}}=(2\pi)^3 2E_{\mathbf{p}}\delta^{(3)}(\mathbf{p}-\mathbf{q}).
\end{equation}

\subsection{Complex Scalar Fields}
Lagrangian:
\begin{equation}
\mathcal{L}=\partial_\mu\phi^\ast\,\partial^\mu\phi - m^2\phi^\ast\phi.
\end{equation}
Canonical momenta:
\begin{equation}
\pi=\frac{\partial\mathcal{L}}{\partial\dot{\phi}}=\dot{\phi}^\ast,
\qquad
\pi^\ast=\frac{\partial\mathcal{L}}{\partial\dot{\phi}^\ast}=\dot{\phi}.
\end{equation}
Mode expansions introduce \textbf{particles and antiparticles}:
\begin{equation}
\phi(x)=\int\frac{d^3p}{(2\pi)^3}\frac{1}{\sqrt{2E_{\mathbf{p}}}}
\left(
a_{\mathbf{p}}e^{-ip\cdot x}+b^\dagger_{\mathbf{p}}e^{ip\cdot x}
\right),
\end{equation}
\begin{equation}
\phi^\ast(x)=\int\frac{d^3p}{(2\pi)^3}\frac{1}{\sqrt{2E_{\mathbf{p}}}}
\left(
a^\dagger_{\mathbf{p}}e^{ip\cdot x}+b_{\mathbf{p}}e^{-ip\cdot x}
\right).
\end{equation}
Nonzero commutators:
\begin{equation}
[a_{\mathbf{p}},a^\dagger_{\mathbf{q}}]=(2\pi)^3\delta^{(3)}(\mathbf{p}-\mathbf{q}),
\qquad
[b_{\mathbf{p}},b^\dagger_{\mathbf{q}}]=(2\pi)^3\delta^{(3)}(\mathbf{p}-\mathbf{q}).
\end{equation}

\paragraph{Conserved charge ($U(1)$ symmetry).}
\begin{equation}
j^\mu=i\left(\phi^\ast\partial^\mu\phi-(\partial^\mu\phi^\ast)\phi\right),
\qquad
Q=\int d^3x\,j^0.
\end{equation}
In terms of operators (schematically):
\begin{equation}
Q=\int\frac{d^3p}{(2\pi)^3}\left(a^\dagger_{\mathbf{p}}a_{\mathbf{p}}-b^\dagger_{\mathbf{p}}b_{\mathbf{p}}\right).
\end{equation}

\subsection{The Heisenberg Picture}
Operators evolve in time, states are time-independent:
\begin{equation}
\mathcal{O}_H(t)=e^{iHt}\mathcal{O}_S e^{-iHt}.
\end{equation}
Equation of motion:
\begin{equation}
\frac{d}{dt}\mathcal{O}_H(t)=i[H,\mathcal{O}_H(t)]+\left(\frac{\partial \mathcal{O}}{\partial t}\right)_H.
\end{equation}
For fields:
\begin{equation}
\dot{\phi}(\mathbf{x},t)=i[H,\phi(\mathbf{x},t)].
\end{equation}

\subsection{Causality and Microcausality}
Relativistic causality requires that measurements at spacelike separation commute:
\begin{equation}
(x-y)^2<0
\quad\Rightarrow\quad
[\phi(x),\phi(y)]=0.
\end{equation}
Define the \textbf{Pauli--Jordan commutator function}:
\begin{equation}
[\phi(x),\phi(y)] = i\Delta(x-y),
\end{equation}
where $\Delta(x)$ vanishes for spacelike $x$.

\subsection{Propagators and Green's Functions}
The \textbf{Feynman propagator} is the time-ordered vacuum correlator:
\begin{equation}
\Delta_F(x-y)
:=\langle 0|T\{\phi(x)\phi(y)\}|0\rangle.
\end{equation}
Time ordering:
\begin{equation}
T\{\phi(x)\phi(y)\}=
\theta(x^0-y^0)\phi(x)\phi(y)+\theta(y^0-x^0)\phi(y)\phi(x).
\end{equation}
Momentum-space representation:
\begin{equation}
\Delta_F(x-y)
=
\int\frac{d^4p}{(2\pi)^4}\,
\frac{i}{p^2-m^2+i\epsilon}\,e^{-ip\cdot(x-y)}.
\end{equation}
It satisfies the Green's function equation:
\begin{equation}
(\square+m^2)\Delta_F(x-y)=-i\delta^{(4)}(x-y).
\end{equation}

\paragraph{Retarded/advanced propagators.}
Retarded Green function:
\begin{equation}
(\square+m^2)\Delta_R(x)=-\delta^{(4)}(x),
\qquad
\Delta_R(x)=0\ \text{for}\ x^0<0.
\end{equation}
Advanced Green function:
\begin{equation}
\Delta_A(x)=0\ \text{for}\ x^0>0.
\end{equation}

\subsection{Applications (What Propagators Are For)}
\begin{itemize}
\item They encode \textbf{free field correlations}:
\[
\langle 0|\phi(x)\phi(y)|0\rangle.
\]
\item They are Green functions used to solve sourced equations:
\[
(\square+m^2)\phi(x)=J(x)
\quad\Rightarrow\quad
\phi(x)=\phi_{\mathrm{hom}}(x)+\int d^4y\,\Delta_R(x-y)J(y).
\]
\item In perturbation theory, $\Delta_F$ becomes the internal line of Feynman diagrams.
\end{itemize}

\subsection{Non-Relativistic Field Theory (Basics)}
Non-relativistic quantum mechanics can be written in second-quantized form using a field
$\psi(\mathbf{x},t)$ with Lagrangian
\begin{equation}
\mathcal{L}
=
i\psi^\dagger\partial_t\psi
-\frac{1}{2m}\nabla\psi^\dagger\cdot\nabla\psi
- V(\mathbf{x})\,\psi^\dagger\psi.
\end{equation}
Canonical momentum:
\begin{equation}
\pi_\psi=\frac{\partial\mathcal{L}}{\partial(\partial_t\psi)}=i\psi^\dagger.
\end{equation}
Quantization (bosons):
\begin{equation}
[\psi(\mathbf{x},t),\psi^\dagger(\mathbf{y},t)]=\delta^{(3)}(\mathbf{x}-\mathbf{y}).
\end{equation}
Hamiltonian:
\begin{equation}
H=\int d^3x\left[\frac{1}{2m}\nabla\psi^\dagger\cdot\nabla\psi
+V(\mathbf{x})\psi^\dagger\psi\right].
\end{equation}
Equation of motion (Heisenberg):
\begin{equation}
i\partial_t\psi(\mathbf{x},t)=\left[-\frac{\nabla^2}{2m}+V(\mathbf{x})\right]\psi(\mathbf{x},t),
\end{equation}
which is the Schr\"odinger equation.

\paragraph{Interactions (typical).}
A common local interaction is
\begin{equation}
\mathcal{L}_{\mathrm{int}}=-\frac{g}{2}(\psi^\dagger\psi)^2.
\end{equation}

\subsection{Summary (Exam Checklist)}
\begin{itemize}
\item Canonical quantization: $[\phi,\pi]=i\delta^{(3)}$ at equal time.
\item Mode expansion introduces $a_{\mathbf{p}},a^\dagger_{\mathbf{p}}$ and particles.
\item Free Hamiltonian is a continuum of oscillators:
\[
H=\int \frac{d^3p}{(2\pi)^3}E_{\mathbf{p}}\left(a^\dagger_{\mathbf{p}}a_{\mathbf{p}}+\frac12\delta^{(3)}(0)\right).
\]
\item Vacuum energy diverges; normal ordering removes the constant shift.
\item Complex scalar $\Rightarrow$ particle/antiparticle operators $a,b$ and conserved $U(1)$ charge.
\item Causality: $[\phi(x),\phi(y)]=0$ for spacelike separation.
\item Feynman propagator:
\[
\Delta_F(x)=\int\frac{d^4p}{(2\pi)^4}\frac{i}{p^2-m^2+i\epsilon}e^{-ip\cdot x}.
\]
\item Non-relativistic fields reproduce Schr\"odinger equation via second quantization.
\end{itemize}
